% Ratios of integers, multidouble's double-double product, and the squares of the relaxed series product.
% \input after the exponential's section of expansion-error-bounds.tex.

\section{Ratios of integers and products of series}\label{sec:recip}

A precision-generic power series needs constants such as $1/k$ in its own rounding: \texttt{DD::ONE / k} widened to
eight parts keeps an error of $2^{-106}$. multidouble's \texttt{recip::ratio\_parts(p, q)} forms $p/q$ as its greedy
expansion by integer arithmetic, and \texttt{mul\_recip(k)} multiplies by it. Its \texttt{series::Relaxed} sums the
Cauchy products of a recurrence in van der Hoeven's squares~\cite{vdH02}, and \texttt{series::mullow} forms products of
series by an exact integer convolution.

\begin{lemma}\label{lem:greedy}
Let $x$ be real, $r_0 = \RN(x)$ and $r_i = \RN(x - r_0 - \dots - r_{i-1})$. Then $|r_{i+1}| \le u|r_i|$, so $r_0,
\dots, r_{n-1}$ lie on $H(1)$, and $|x - \sum_{i<n} r_i| \le u^n(1 + u)|x|$ for $n \ge 1$. The $\TwoSum$ of $r_0$ and
$r_1$ is a double word of their sum, and no double word $(h, \ell)$, $h = \RN(h + \ell)$, is nearer $x$.
\end{lemma}

\begin{proof}
What $r_i$ leaves is at most $u|r_i|$, a float, so its rounding $r_{i+1}$ is too; and the remainder after $n$ parts
is at most $u|r_{n-1}| \le u^n|r_0| \le u^n(1+u)|x|$. Let $t = x - r_0$. A double word with $h = r_0$ is no nearer,
as $\ell$ is a float and $r_1 = \RN(t)$. Otherwise, say $h > r_0$: $x$ lies at or below the midpoint of $r_0$ and $h$,
and $h + \ell$ at or above it, so $|x - h - \ell| \ge (h - r_0)/2 - t$. The next float $s > r_0$ is at most $h$, and
$g = (s - r_0)/2$ is a float, half the spacing of the floats at $r_0$, with $t \le g$; so $|t - r_1| \le g - t \le
|x - h - \ell|$, and for $t \le 0$ already $|t - r_1| \le |t|$. The case $h < r_0$ is its mirror.
\end{proof}

The pair $(r_0, r_1)$ itself need not be a double word: where $r_1$ is half an ulp of an odd $r_0$, $r_0 + r_1$ is a
tie that rounds to the even neighbor. This needs $q \ge 2^{53}$ for $x = p/q$, and occurs for $q = 191516586041457763$;
the $\TwoSum$ keeps the value and moves $r_0$.

\begin{proposition}\label{prop:ratio}
For integers $|p| \le 2^{63}$ and $1 \le q < 2^{60}$, \texttt{ratio\_parts(p, q)} returns the first eight greedy parts
of $p/q$. Each step keeps what is left as $(-1)^\nu n/(q2^e)$ with $n$ an integer, puts $n2^s/q$ in $[2^{52}, 2^{53})$
by a shift $s \in [-11, 112]$, rounds it to an integer, to nearest with ties to even, by one division, and scales it by
$2^{-(e + s)}$. Its dividends stay below $2^{127}$, its divisors below $2^{72}$ and its products below $2^{127}$, so
u128 and i128 hold them, and every scale is a normal binary64 number. So \texttt{DD::recip\_int(k)} and
\texttt{DD::from\_ratio(p, q)}, the $\TwoSum$ of the first two parts, are the double words nearest $1/k$ and $p/q$,
within $u^2(1 + u)$, and the quad- and octo-double constants, the first four and eight parts, lie on $H(1)$ within
$u^4(1 + u)$ and $u^8(1 + u)$, relatively.
\end{proposition}

\begin{proof}
With $s_0 = 52 - (\mathrm{bits}(n) - \mathrm{bits}(q))$, $n2^{s_0}/q$ lies in $(2^{51}, 2^{53})$, and the shift
is $s_0$ or $s_0 + 1$. The rounding of the quotient $q_0$ up where the remainder exceeds half the divisor, or equals
it with $q_0$ odd, is $\RN$ of $n2^s/q$, and the part is $\RN$ of what is left; the new remainder is what the part
leaves, at most half the divisor. So $n \le 2^{63}$ before the first part, $n < q2^{11}$ before the second and $n \le
q/2$ before the others, which gives $s \ge -11$, $s \ge 41$ and $s \ge 52$, and $e \le 112i$ before part $i$, so $e +
s \le 896$ for the eight parts.
\end{proof}

multidouble's double-double product is not one of the published algorithms~\cite{JMP17}: it adds both cross
products to the error of $\TwoProd$ by fused multiply-adds.

\begin{proposition}\label{prop:ddmul}
For double words $x$ and $y$ and $p \ge 3$, let $(c_h, c_\ell) = \TwoProd(x_h, y_h)$, $t_1 = \RN(x_hy_\ell + c_\ell)$,
$t_2 = \RN(x_\ell y_h + t_1)$ and $(z_h, z_\ell) = \FastTwoSum(c_h, t_2)$. Then $z$ is a double word with $|z - xy|
\le (6 + 2u)u^2/(1 - u)^2\,|xy|$, at most $6.000001\,u^2|xy|$ at $p = 53$.
\end{proposition}

\begin{proof}
Write $P = |x_hy_h|$. The error is $(t_2 - (x_\ell y_h + t_1)) + (t_1 - (x_hy_\ell + c_\ell)) - x_\ell y_\ell$. With
$|x_\ell| \le u|x_h|$, $|y_\ell| \le u|y_h|$ and $|c_\ell| \le uP$, $|x_hy_\ell + c_\ell| \le 2uP$ and $|x_\ell y_h +
t_1| \le (3 + 2u)uP$, so the error is at most $(6 + 2u)u^2P$; and $|t_2| \le |c_h|$, so $\FastTwoSum$ is exact. As
$|xy| \ge (1 - u)^2P$, the bound follows.
\end{proof}

multidouble's tests check the threshold $5u^2$; the bound here holds for every pair of double words.

\begin{corollary}\label{cor:mulrecip}
For a double word $x$, \texttt{DD::mul\_recip(k)} is within $(\varepsilon_{\rm P}(1 + u^2(1 + u)) + u^2(1 + u))|x/k|$
of $x/k$, $\varepsilon_{\rm P} = (6 + 2u)u^2/(1 - u)^2$, at most $7.000001\,u^2|x/k|$ at $p = 53$. For $x$ on $H(2)$,
\texttt{OD::mul\_recip(k)} is within $(\varepsilon_8(1 + u^8(1 + u)) + u^8(1 + u))|x/k|$, $\varepsilon_8 = u^8 +
\kappa u^9((1 - 2u)/(1 - 4u))^2$ with Table~\ref{tab:kappa}'s $\kappa \le 2^{37.93}$, at most $2.00003\,u^8|x/k|$.
\end{corollary}

\begin{proof}
A product within $\varepsilon$ of $xc$, by a constant within $\delta$ of $1/k$, is within $\varepsilon(1 + \delta) +
\delta$ of $x/k$, relatively; $\varepsilon$ is Proposition~\ref{prop:ddmul}'s and Theorem~\ref{thm:t3}'s for MUL on
$H(2)$, $\delta$ Lemma~\ref{lem:greedy}'s.
\end{proof}

\begin{proposition}\label{prop:relaxed}
Let $s_0 = 2^a$ be the least side. At step $e$, \texttt{Relaxed} forms for each side $s = 2^j \ge s_0$ with $(m + 1)s =
e - 1$, $m \ge 1$, the square of the terms $x_iy_j$ with $i \in [ms + 1, (m + 1)s]$ and $j \in [s + 1, 2s]$, and for $m
\ge 2$ its transpose. Each term with $i, j \ge s_0 + 1$ lies in exactly one square of one step, which takes
coefficients below $e$ and gives orders $e + 1$ and above. At order $k \ge 2d$, $d = s_0 + 1$, the terms with $i < d$
or $k - i < d$ are summed directly, once each.
\end{proposition}

\begin{proof}
With $i' = i - 1$ and $j' = j - 1$, the term lies in the square of side $2^{\lfloor\log_2 j'\rfloor}$ with $m =
\lfloor i'/s\rfloor$ where $\lfloor\log_2 i'\rfloor \ge \lfloor\log_2 j'\rfloor$, and otherwise in the transpose of
side $2^{\lfloor\log_2 i'\rfloor}$, with $m \ge 2$; a square's $j' \in [s, 2s)$ fixes $s$, and $i' \in [ms, (m+1)s)$
fixes $m$, so no other square holds it. The square of $(s, m)$ belongs to step $(m + 1)s + 1$ alone.
\end{proof}

\texttt{series::mullow} multiplies two series block by block, as Arb does~\cite{Johansson17} after van der
Hoeven~\cite{vdH08}. In a block of spread $h$ it scales the coefficient $a_i$ to $\alpha_i = 2^{k_i}a_i$, $W - h \le
e_i + k_i \le W$, where $2^{e_i - 1} \le |a_{i,0}| < 2^{e_i}$ and $W = 16t - 2$; rounds each part to an integer; and
writes their sum as $t$ balanced 16-bit digits, $|d_r| \le 2^{15}$. Binary64 FFTs convolve the digit sequences, and
each weight $C_q = \sum_{r + r' = q} d_rd'_{r'}$, $q \ge q_0$, is rounded to an integer; the weights below $q_0$ are
not formed. The integer of an output index, carried into balanced digits, is cut into eight 53-bit parts from its top
bit and deposited in the index's accumulator.

\begin{proposition}\label{prop:mullow}
Let $p = 53$, let each part of a coefficient be at most $u$ times the one before, and let the rounded weights of the
inverse transform be the exact $C_q$.
\begin{itemize}
\item[(a)] Every floating-point operation that forms the digits is exact but $\RN(x + M)$, $M = 3 \cdot 2^{51}$, which
rounds a scaled remainder $x$ to an integer. A zero coefficient has zero digits; any other's digits are an integer
within $1/(1 - u)$ of $\alpha_i$, and $2^{W-h-1}(1 - 2u) \le |\alpha_i| \le 2^W(1 + 2u)$.
\item[(b)] Let $W \ge \mathrm{prec} + h + 5$ for both blocks, $l_at \le 2^{20}$, and $q_0 = 0$ or $16q_0 \le 2W - h_a
- h_b - \mathrm{prec} - 16 - \lambda$ with $t\min(l_a, l_b) \le 2^\lambda$. A block pair's integer at index $m$, times
its scale, is within $1.13 \cdot 2^{-\mathrm{prec}}\sum|a_ib_j|$ of $\sum a_ib_j$, both sums over the pair's terms
with $i + j = m$.
\item[(c)] At octo-double, $\mathrm{prec} = 424$, its eight parts lie on $H(2)$ within $2^{-423}$ of it. An index whose
accumulator receives $K \le 50$ terms, block pairs and products formed term by term, gets $c_m$ within $2^{-424}|c_m|
+ 4.15 \cdot 2^{-424}\sum_{i+j=m}|a_ib_j|$.
\end{itemize}
\end{proposition}

The octo-double bound is that of a dot product, at most $5.15 \cdot 2^{-424}\sum|a_ib_j|$. The transform's exactness
is assumed. The code checks that each weight lies within $1/4$ of an integer, which a transform error near 1 also
passes; a bound below $1/2$ on the binary64 FFT's error would remove the assumption. At double- and quad-double (a) and
(b) hold; their read-out and accumulators are not proved.

\begin{proof}
(a) $\RN(\RN(x + M) - M)$ is an integer within $1/2$ of $x$ for $|x| < 2^{51}$: $x + M$ lies
in $(2^{52}, 2^{53})$, where the floats are the integers. The code's $r_{\rm hi}$ keeps each argument of this rounding
below $2^{16}$. The update $\mathit{rem} - q2^{16r}$ of the remainder is exact by Sterbenz's lemma, and so are the
digits' sums and carries, integers below $2^{53}$; so the loop computes its model in real numbers. A part loses at most
itself below 1. A part of at least 1 is a multiple of its lowest digit's weight by the code's $r_{\rm lo}$, and loses
nothing, unless $r_{\rm lo} = 0$, where it loses at most $1/2$. A part of at least $2^{52}$ is an integer, so only the
last part of at least 1 can lose $1/2$, and the losses sum to at most $1/(1 - u)$. The top digit carries nothing out,
as $|\alpha_i| + 2 \le 0.49 \cdot 2^{16t}$.

(b) A term errs by at most $(\eta/L_a + \eta/L_b + \eta^2/(L_aL_b))|\alpha_i\beta_j|$ with $\eta = 1/(1 - u)$ and $L =
2^{W-h-1}(1 - 2u)$. Each weight is at most $t2^{30}$ per term, so the dropped weights add at most $t2^{30}2^{16q_0}/
((2^{16} - 1)L_aL_b)$ times $|\alpha_i\beta_j|$, and the four terms sum to at most $1.13 \cdot 2^{-\mathrm{prec}}$.
The weights are below $2^{50}$, so their carry into digits is exact.

(c) With $h_b$ the top bit of the integer, the first $j + 1$ chunks are the integer truncated below bit $b_j = h_b - 52
- 53j$, so the eight lose less than $2^{h_b - 423}$. Chunk $j \ge 1$ is below $2^{53 + b_j}$, at most $2^{1 - 53j} \le
(2u)^j$ times the first. Theorem~\ref{thm:acc} for $K \le 50$ terms on $H(2)$ (Table~\ref{tab:acc-k}) gives
$2^{-424}(|Z| + \sum_k\mu_k)$, and $\sum_k\mu_k \le 1.0001\sum|a_ib_j|$.
\end{proof}
